\chapter{Boolean Analysis — Canonical D Tree}
%%% generated by the blueprint pipeline from TCSlib.BooleanAnalysis.Switching.CanonicalDTree
\section{Overview}

This module builds the \emph{canonical decision tree} for a restricted DNF $f|_\rho$,
following Razborov's construction: at each stage one selects the first term not killed
by the current restriction $\rho$, branches on all of that term's free variables, and
recurses. The bulk of the file establishes that this tree correctly computes $f|_\rho$,
is independent of the recursion fuel, and that its depth bounds the decision-tree depth
of $f|_\rho$.

\section{Declarations}

\begin{lemma}[Fixing a free variable lowers the free count]
\lean{SwitchingLemma2.numFree_update_lt}
\label{SwitchingLemma2.numFree_update_lt}
\leanok
\uses{SwitchingLemma2.Restriction, SwitchingLemma2.Restriction.freeVars, SwitchingLemma2.Restriction.numFree}
\difficulty{4}
Let $\rho$ be a restriction on $n$ variables, and let $v$ be a free variable of $\rho$.
For any Boolean value $b$, let $\rho'$ be the restriction that agrees with $\rho$
everywhere except at $v$, where it assigns the fixed value $b$. Then the number of free
variables of $\rho'$ is strictly less than the number of free variables of $\rho$.
\end{lemma}

\begin{lemma}[Existence of a free variable when no term is fixed and not all are killed]
\lean{SwitchingLemma2.exists_free_of_not_killed_not_fixed}
\statementsource{switching-lemma}{pdf-b5e074215b9e-p002-b002}
\label{SwitchingLemma2.exists_free_of_not_killed_not_fixed}
\leanok
\uses{DNF, SwitchingLemma2.Literal.killedBy, SwitchingLemma2.Restriction, SwitchingLemma2.Restriction.freeVars, SwitchingLemma2.Term.fixedBy, SwitchingLemma2.Term.killedBy}
\difficulty{5}
Let $f$ be a DNF formula on $n$ variables and let $\rho$ be a restriction on those
variables. Suppose that not every term of $f$ is killed by $\rho$, and that no term of
$f$ is fixed by $\rho$. Then $\rho$ leaves at least one variable free; that is, there
exists an index $v \in \mathrm{Fin}\,n$ with $\rho(v)$ unset.
\end{lemma}

\begin{definition}[Branch variable selection]
\lean{SwitchingLemma2.selectBranchVar}
\label{SwitchingLemma2.selectBranchVar}
\leanok
\statementsource{switching-lemma}{pdf-b5e074215b9e-p002-b002}
\uses{DNF, SwitchingLemma2.Restriction, SwitchingLemma2.Restriction.freeVars, SwitchingLemma2.Term.killedBy}
Given a DNF $f$ and restriction $\rho$, returns an optional variable in $\mathrm{Fin}\,n$:
the variable on which the canonical tree should next branch, drawn from the free variables
of the first term not killed by $\rho$.
\end{definition}

\begin{lemma}[Branch-variable selection returns a free variable]
\lean{SwitchingLemma2.selectBranchVar_spec}
\label{SwitchingLemma2.selectBranchVar_spec}
\leanok
\proofsource{switching-lemma}{pdf-b5e074215b9e-p002-b002}
\proofstep
  {BoolCircuit.Lit}
  {context}
  {switching-lemma}
  {pdf-b5e074215b9e-p001-b004}
\proofstep
  {BoolCircuit.Lit.eval}
  {context}
  {switching-lemma}
  {pdf-b5e074215b9e-p001-b004}
\proofstep
  {Literal}
  {context}
  {switching-lemma}
  {pdf-b5e074215b9e-p001-b004}
\proofstep
  {Term}
  {context}
  {switching-lemma}
  {pdf-b5e074215b9e-p001-b004}
\proofstep
  {DNF}
  {context}
  {switching-lemma}
  {pdf-b5e074215b9e-p001-b004}
\proofstep
  {SwitchingLemma2.Restriction}
  {context}
  {switching-lemma}
  {pdf-b5e074215b9e-p001-b007}
\proofstep
  {SwitchingLemma2.Restriction.freeVars}
  {context}
  {switching-lemma}
  {pdf-b5e074215b9e-p001-b007}
\proofstep
  {SwitchingLemma2.Literal.fixedBy}
  {context}
  {switching-lemma}
  {pdf-b5e074215b9e-p002-b004}
\proofstep
  {SwitchingLemma2.Term.killedBy}
  {context}
  {switching-lemma}
  {pdf-b5e074215b9e-p002-b002}
\proofstep
  {SwitchingLemma2.Term.fixedBy}
  {context}
  {switching-lemma}
  {pdf-b5e074215b9e-p002-b002}
\proofstep
  {SwitchingLemma2.selectBranchVar}
  {context}
  {switching-lemma}
  {pdf-b5e074215b9e-p002-b002}
\uses{DNF, SwitchingLemma2.Literal.killedBy, SwitchingLemma2.Restriction, SwitchingLemma2.Restriction.freeVars, SwitchingLemma2.Term.fixedBy, SwitchingLemma2.Term.killedBy, SwitchingLemma2.canonicalDTree, SwitchingLemma2.selectBranchVar}
\difficulty{5}
Let $f$ be a DNF formula on $n$ variables and $\rho$ a restriction. Suppose that at
least one term of $f$ is not killed by $\rho$, and that no term of $f$ is fixed by
$\rho$. Then the branch-variable selection procedure applied to $f$ and $\rho$ succeeds,
returning a variable $v$ that is free in $\rho$.
\end{lemma}

\begin{definition}[Sub-tree for a term]
\lean{SwitchingLemma2.termSubTree}
\label{SwitchingLemma2.termSubTree}
\leanok
\statementsource{switching-lemma}{
  pdf-b5e074215b9e-p002-b002,
  pdf-b5e074215b9e-p002-b003
}
\uses{DecisionTree, Literal, SwitchingLemma2.Restriction, SwitchingLemma2.Restriction.freeVars}
Builds a complete sub-tree for a term, queried as a list of literals: it branches on each
free variable of the term in order, and at every leaf invokes a continuation
$\mathrm{cont} : \mathrm{Restriction}\,n \to \mathrm{DecisionTree}\,n$ with the restriction updated
along that root-to-leaf path. Non-free literals are skipped.
\end{definition}

\begin{definition}[Canonical decision tree]
\lean{SwitchingLemma2.canonicalDTree}
\label{SwitchingLemma2.canonicalDTree}
\leanok
\statementsource{switching-lemma}{
  pdf-b5e074215b9e-p002-b002,
  pdf-b5e074215b9e-p002-b003
}
\uses{DNF, DecisionTree, SwitchingLemma2.Restriction, SwitchingLemma2.Restriction.numFree}
The canonical decision tree for $f|_\rho$, following Razborov's construction. It is defined
through a fuel-driven helper $\mathrm{canonicalDTree.go}$ initialized with fuel
$\rho.\mathrm{numFree} + \text{(positive)}$, repeatedly selecting the first non-killed term and
expanding its free variables.
\end{definition}

\begin{lemma}[Updating a restriction at a free variable to match the point]
\lean{SwitchingLemma2.extend_update_self}
\label{SwitchingLemma2.extend_update_self}
\leanok
\uses{SwitchingLemma2.Restriction, SwitchingLemma2.Restriction.extend, SwitchingLemma2.Restriction.freeVars}
\difficulty{3}
Let $\rho$ be a restriction on $n$ variables, let $x : \mathrm{Fin}\,n \to
\mathrm{Bool}$ be a point, and let $v \in \mathrm{Fin}\,n$ be a free variable of $\rho$
with $x(v) = b$. Let $\rho[v \mapsto b]$ be the restriction obtained from $\rho$ by
fixing the value of variable $v$ to $b$ (and leaving every other variable as in $\rho$).
Then $\rho[v \mapsto b]$ and $\rho$ have the same extension along $x$; that is, $(\rho[v
\mapsto b]).\mathrm{extend}\,x$ and $\rho.\mathrm{extend}\,x$ agree as Boolean
assignments.
\end{lemma}

\begin{lemma}[Semantic correctness of the term sub-tree]
\lean{SwitchingLemma2.termSubTree_eval}
\label{SwitchingLemma2.termSubTree_eval}
\leanok
\uses{DecisionTree, DecisionTree.eval, Literal, SwitchingLemma2.Restriction, SwitchingLemma2.Restriction.freeVars, SwitchingLemma2.termSubTree}
\difficulty{4}
Fix $n$, and let $\ell_1,\dots,\ell_m$ be a list of literals over $n$ variables, $\rho$
a restriction on those variables, $\mathrm{cont}$ a map sending each restriction to a
decision tree, and $x \in \{\mathrm{true},\mathrm{false}\}^n$ an assignment. Then
evaluating the term sub-tree built from $\ell_1,\dots,\ell_m$, $\rho$, and
$\mathrm{cont}$ at $x$ gives the same Boolean value as evaluating
$\mathrm{cont}(\rho^\ast)$ at $x$, where $\rho^\ast$ is obtained from $\rho$ by
processing the literals in order: for each $\ell_i$, if its variable is free in the
current restriction, that variable is fixed to $x$'s value there, and otherwise the
restriction is left unchanged.
\end{lemma}

\begin{lemma}[Fixing free variables to a point preserves its extension]
\lean{SwitchingLemma2.termSubTree_extend_eq}
\label{SwitchingLemma2.termSubTree_extend_eq}
\leanok
\uses{Literal, SwitchingLemma2.Restriction, SwitchingLemma2.Restriction.extend, SwitchingLemma2.Restriction.freeVars, SwitchingLemma2.extend_update_self}
\difficulty{3}
Let $\rho$ be a restriction on $n$ variables, let $x : \mathrm{Fin}\,n \to
\mathrm{Bool}$ be a point, and let $\ell_1, \dots, \ell_m$ be a list of literals over
these variables. Form a sequence of restrictions $\rho_0 = \rho, \rho_1, \dots, \rho_m$
by processing the literals in order: for each $k$, writing $v_k$ for the variable of
$\ell_k$, if $v_k$ is free in $\rho_{k-1}$ then let $\rho_k$ agree with $\rho_{k-1}$
except that $v_k$ is fixed to the value $x(v_k)$, and otherwise let $\rho_k =
\rho_{k-1}$. Then the extension of $\rho_m$ along $x$ equals the extension of $\rho$
along $x$.
\end{lemma}

\begin{lemma}[Fixed variables survive the literal fold]
\lean{SwitchingLemma2.termSubTree_foldl_preserves_nonnone}
\label{SwitchingLemma2.termSubTree_foldl_preserves_nonnone}
\leanok
\uses{Literal, SwitchingLemma2.Restriction, SwitchingLemma2.Restriction.freeVars}
\difficulty{3}
Let $\rho : \mathrm{Fin}\,n \to \mathrm{Option}\,\mathrm{Bool}$ be a restriction on $n$
variables, let $x : \mathrm{Fin}\,n \to \mathrm{Bool}$ be a Boolean assignment, and let
$\ell_1, \dots, \ell_m$ be literals on these variables. Process the literals from left
to right, starting from $\rho$: at each literal, if its variable is still free in the
restriction accumulated so far, fix that variable to the value $x$ assigns it, and
otherwise change nothing. Then any variable $v$ that already carries a fixed value under
$\rho$ (that is, $\rho(v)$ leaves $v$ set rather than free) still carries a fixed value
in the final restriction; the fold never frees an already-fixed variable.
\end{lemma}

\begin{lemma}[Folding literals fixes each initially free variable]
\lean{SwitchingLemma2.termSubTree_foldl_sets_member}
\label{SwitchingLemma2.termSubTree_foldl_sets_member}
\leanok
\uses{Literal, SwitchingLemma2.Restriction, SwitchingLemma2.Restriction.freeVars, SwitchingLemma2.termSubTree_foldl_preserves_nonnone}
\difficulty{4}
Fix $n$ variables, let $x : \mathrm{Fin}\,n \to \mathrm{Bool}$ be a Boolean assignment,
let $\rho$ be a restriction on these variables, and let $\ell_1, \dots, \ell_m$ be a
list of literals. Process the literals from left to right, maintaining a current
restriction that begins as $\rho$: at each literal $\ell$, if its underlying variable
$\ell.\mathrm{var}$ is free in the current restriction, update that restriction by
fixing $\ell.\mathrm{var}$ to the value $x(\ell.\mathrm{var})$, and otherwise leave the
restriction unchanged; write $\rho^*$ for the restriction obtained after the whole list
has been processed. Then for every literal $\ell$ in the list whose variable is free in
the original restriction $\rho$, the value $\rho^*(\ell.\mathrm{var})$ is set, i.e.\
$\ell.\mathrm{var}$ is not free in $\rho^*$.
\end{lemma}

\begin{lemma}[Sequential fixing strictly reduces the free count]
\lean{SwitchingLemma2.termSubTree_foldl_numFree_lt}
\label{SwitchingLemma2.termSubTree_foldl_numFree_lt}
\leanok
\uses{Literal, SwitchingLemma2.Restriction, SwitchingLemma2.Restriction.freeVars, SwitchingLemma2.Restriction.numFree, SwitchingLemma2.termSubTree_foldl_preserves_nonnone, SwitchingLemma2.termSubTree_foldl_sets_member}
\difficulty{4}
Let $\rho$ be a restriction on $n$ variables, let $x : \mathrm{Fin}\,n \to
\mathrm{Bool}$ be a Boolean assignment, and let $\ell_1, \dots, \ell_k$ be a list of
literals on $n$ variables. Process the literals from left to right, starting from
$\rho$: at each step, if the current literal's variable is still free in the restriction
obtained so far, fix that variable to the value $x$ assigns it, and otherwise leave the
restriction unchanged; write $\rho'$ for the restriction produced after all $k$ steps.
If at least one literal in the list has its variable free in the original restriction
$\rho$, then $\rho'$ has strictly fewer free variables than $\rho$.
\end{lemma}

\begin{lemma}[Correctness of the fuel-bounded decision-tree construction]
\lean{SwitchingLemma2.canonicalDTree_go_correct}
\label{SwitchingLemma2.canonicalDTree_go_correct}
\leanok
\proofsource{switching-lemma}{pdf-b5e074215b9e-p002-b002}
\uses{DNF, DNF.eval, DecisionTree.eval, SwitchingLemma2.Literal.fixedBy_eval_true, SwitchingLemma2.Literal.killedBy, SwitchingLemma2.Literal.killedBy_eval_false, SwitchingLemma2.Restriction, SwitchingLemma2.Restriction.extend, SwitchingLemma2.Restriction.freeVars, SwitchingLemma2.Restriction.numFree, SwitchingLemma2.Term.fixedBy, SwitchingLemma2.Term.killedBy, SwitchingLemma2.list_all_eq_false_of_mem, SwitchingLemma2.list_any_eq_false, SwitchingLemma2.restrictFn, SwitchingLemma2.termSubTree_eval, SwitchingLemma2.termSubTree_extend_eq, SwitchingLemma2.termSubTree_foldl_numFree_lt, Term.eval}
\difficulty{6}
Let $f$ be a DNF formula on $n$ Boolean variables, let $\rho$ be a restriction of those
variables, and let $m$ be a natural number strictly greater than the number of free
variables of $\rho$. Denote by $T$ the decision tree obtained from $f$ and $\rho$ by the
canonical fuel-bounded recursion run with fuel budget $m$. Then $T$ computes the
restriction of $f$ by $\rho$: for every assignment $x : \mathrm{Fin}\,n \to
\mathrm{Bool}$, evaluating $T$ at $x$ returns $f$ applied to the extension of $x$ along
$\rho$ — the assignment that uses $\rho$'s fixed value on each constrained coordinate
and $x$'s value on each free coordinate.
\end{lemma}

\begin{lemma}[Correctness of the canonical decision tree]
\lean{SwitchingLemma2.canonicalDTree_correct}
\label{SwitchingLemma2.canonicalDTree_correct}
\leanok
\proofsource{switching-lemma}{pdf-b5e074215b9e-p002-b002}
\proofstep
  {BoolCircuit.Lit}
  {context}
  {switching-lemma}
  {pdf-b5e074215b9e-p001-b004}
\proofstep
  {BoolCircuit.Lit.eval}
  {context}
  {switching-lemma}
  {pdf-b5e074215b9e-p001-b004}
\proofstep
  {Literal}
  {context}
  {switching-lemma}
  {pdf-b5e074215b9e-p001-b004}
\proofstep
  {DecisionTree.eval}
  {context}
  {switching-lemma}
  {pdf-b5e074215b9e-p001-b005}
\proofstep
  {Term}
  {context}
  {switching-lemma}
  {pdf-b5e074215b9e-p001-b004}
\proofstep
  {DNF}
  {context}
  {switching-lemma}
  {pdf-b5e074215b9e-p001-b004}
\proofstep
  {DecisionTree}
  {context}
  {switching-lemma}
  {pdf-b5e074215b9e-p001-b005}
\proofstep
  {SwitchingLemma2.Restriction}
  {context}
  {switching-lemma}
  {pdf-b5e074215b9e-p001-b007}
\proofstep
  {SwitchingLemma2.Restriction.freeVars}
  {context}
  {switching-lemma}
  {pdf-b5e074215b9e-p001-b007}
\proofstep
  {SwitchingLemma2.Restriction.numFree}
  {context}
  {switching-lemma}
  {pdf-b5e074215b9e-p001-b007}
\proofstep
  {SwitchingLemma2.Restriction.extend}
  {context}
  {switching-lemma}
  {pdf-b5e074215b9e-p001-b007}
\proofstep
  {SwitchingLemma2.Literal.fixedBy}
  {context}
  {switching-lemma}
  {pdf-b5e074215b9e-p002-b004}
\proofstep
  {SwitchingLemma2.Term.killedBy}
  {context}
  {switching-lemma}
  {pdf-b5e074215b9e-p002-b002}
\proofstep
  {SwitchingLemma2.Term.fixedBy}
  {context}
  {switching-lemma}
  {pdf-b5e074215b9e-p002-b002}
\proofstep
  {SwitchingLemma2.restrictFn}
  {context}
  {switching-lemma}
  {pdf-b5e074215b9e-p001-b007}
\proofstep
  {SwitchingLemma2.Literal.killedBy_eval_false}
  {context}
  {switching-lemma}
  {pdf-b5e074215b9e-p002-b002}
\proofstep
  {SwitchingLemma2.Literal.fixedBy_eval_true}
  {context}
  {switching-lemma}
  {pdf-b5e074215b9e-p002-b002}
\proofstep
  {SwitchingLemma2.list_any_eq_false}
  {context}
  {switching-lemma}
  {pdf-b5e074215b9e-p002-b002}
\proofstep
  {SwitchingLemma2.list_all_eq_false_of_mem}
  {context}
  {switching-lemma}
  {pdf-b5e074215b9e-p002-b002}
\proofstep
  {SwitchingLemma2.termSubTree}
  {context}
  {switching-lemma}
  {pdf-b5e074215b9e-p002-b002}
\proofstep
  {SwitchingLemma2.canonicalDTree}
  {context}
  {switching-lemma}
  {pdf-b5e074215b9e-p002-b002}
\proofstep
  {SwitchingLemma2.extend_update_self}
  {context}
  {switching-lemma}
  {pdf-b5e074215b9e-p002-b002}
\proofstep
  {SwitchingLemma2.termSubTree_eval}
  {context}
  {switching-lemma}
  {pdf-b5e074215b9e-p002-b002}
\proofstep
  {SwitchingLemma2.termSubTree_extend_eq}
  {context}
  {switching-lemma}
  {pdf-b5e074215b9e-p002-b002}
\proofstep
  {SwitchingLemma2.termSubTree_foldl_preserves_nonnone}
  {context}
  {switching-lemma}
  {pdf-b5e074215b9e-p002-b002}
\proofstep
  {SwitchingLemma2.termSubTree_foldl_sets_member}
  {context}
  {switching-lemma}
  {pdf-b5e074215b9e-p002-b002}
\proofstep
  {SwitchingLemma2.termSubTree_foldl_numFree_lt}
  {context}
  {switching-lemma}
  {pdf-b5e074215b9e-p002-b002}
\proofstep
  {SwitchingLemma2.canonicalDTree_go_correct}
  {direct}
  {switching-lemma}
  {pdf-b5e074215b9e-p002-b002}
\uses{DNF, DNF.eval, DecisionTree.eval, SwitchingLemma2.Restriction, SwitchingLemma2.canonicalDTree, SwitchingLemma2.canonicalDTree_go_correct, SwitchingLemma2.restrictFn}
\difficulty{1}
Let $f$ be a DNF formula on $n$ Boolean variables and let $\rho$ be a restriction of
those variables. Then the canonical decision tree of $f$ under $\rho$ computes exactly
the restriction of $f$ by $\rho$: for every Boolean assignment $x : \mathrm{Fin}\,n \to
\mathrm{Bool}$, evaluating the tree at $x$ yields the value of $f$ on the assignment
that takes $\rho(i)$ wherever $\rho$ fixes coordinate $i$ and $x_i$ on the free
coordinates.
\end{lemma}

\begin{lemma}[Fixing a variable does not increase the free count]
\lean{SwitchingLemma2.numFree_update_le}
\label{SwitchingLemma2.numFree_update_le}
\leanok
\uses{SwitchingLemma2.Restriction, SwitchingLemma2.Restriction.freeVars, SwitchingLemma2.Restriction.numFree}
\difficulty{3}
Let $\rho$ be a restriction on $n$ variables, let $v \in \mathrm{Fin}\,n$, and let $b$
be a Boolean value. Write $\rho[v \mapsto b]$ for the restriction that agrees with
$\rho$ everywhere except at $v$, where it takes the fixed value $b$. Then the number of
free variables of $\rho[v \mapsto b]$ is at most the number of free variables of $\rho$.
\end{lemma}

\begin{lemma}[Continuation-extensionality of the term sub-tree]
\lean{SwitchingLemma2.termSubTree_cont_congr}
\label{SwitchingLemma2.termSubTree_cont_congr}
\leanok
\uses{DecisionTree, Literal, SwitchingLemma2.Restriction, SwitchingLemma2.Restriction.freeVars, SwitchingLemma2.Restriction.numFree, SwitchingLemma2.numFree_update_le, SwitchingLemma2.termSubTree}
\difficulty{4}
Let $\ell$ be a list of literals on $n$ variables and let $\rho$ be a restriction. For a
continuation $c : \mathrm{Restriction}\,n \to \mathrm{DecisionTree}\,n$, write $T(\ell,
\rho, c)$ for the term sub-tree built from $\ell$ starting at $\rho$ with continuation
$c$. If two continuations $c_1$ and $c_2$ agree on every restriction $\rho'$ whose
number of free variables is at most the number of free variables of $\rho$ — that is,
$c_1(\rho') = c_2(\rho')$ whenever $\rho'$ has no more free variables than $\rho$ — then
$T(\ell, \rho, c_1) = T(\ell, \rho, c_2)$.
\end{lemma}

\begin{lemma}[Strict continuation-invariance of the term sub-tree]
\lean{SwitchingLemma2.termSubTree_cont_congr_strict}
\label{SwitchingLemma2.termSubTree_cont_congr_strict}
\leanok
\uses{DecisionTree, Literal, SwitchingLemma2.Restriction, SwitchingLemma2.Restriction.freeVars, SwitchingLemma2.Restriction.numFree, SwitchingLemma2.numFree_update_lt, SwitchingLemma2.termSubTree, SwitchingLemma2.termSubTree_cont_congr}
\difficulty{4}
Let $L$ be a list of literals on $n$ variables and $\rho$ a restriction on those
variables, and suppose at least one literal in $L$ has its variable free under $\rho$.
Let $c_1$ and $c_2$ be two continuations — maps sending each restriction to a decision
tree — that agree on every restriction having strictly fewer free variables than $\rho$.
Then the two term sub-trees built from $L$ and $\rho$, one using the continuation $c_1$
and the other using $c_2$, are identical.
\end{lemma}

\begin{lemma}[Fuel invariance of the canonical decision-tree recursion]
\lean{SwitchingLemma2.canonicalDTree_go_fuel_invariant}
\label{SwitchingLemma2.canonicalDTree_go_fuel_invariant}
\leanok
\uses{DNF, SwitchingLemma2.Literal.killedBy, SwitchingLemma2.Restriction, SwitchingLemma2.Restriction.freeVars, SwitchingLemma2.Restriction.numFree, SwitchingLemma2.Term.fixedBy, SwitchingLemma2.Term.killedBy, SwitchingLemma2.termSubTree_cont_congr_strict}
\difficulty{6}
Let $f$ be a DNF formula on $n$ variables, and let $\rho$ be a restriction on those
variables having $k$ free variables. Write $G(f, m, \rho)$ for the fuel-bounded
recursion that builds a decision tree from the DNF $f$ and the restriction $\rho$, in
which the natural-number fuel budget $m$ serves as an upper bound on the recursion
depth. Then for any two fuel budgets $m_1$ and $m_2$ that both strictly exceed $k$, the
two trees coincide: $G(f, m_1, \rho) = G(f, m_2, \rho)$.
\end{lemma}

\begin{lemma}[Fuel invariance of the canonical decision tree continuation]
\lean{SwitchingLemma2.cont_eq_canonicalDTree}
\label{SwitchingLemma2.cont_eq_canonicalDTree}
\leanok
\uses{DNF, SwitchingLemma2.Literal.fixedBy, SwitchingLemma2.Literal.killedBy, SwitchingLemma2.Restriction, SwitchingLemma2.Restriction.numFree, SwitchingLemma2.Term.fixedBy, SwitchingLemma2.canonicalDTree, SwitchingLemma2.canonicalDTree_go_fuel_invariant, Term}
\difficulty{4}
The canonical decision tree of a DNF formula $f$ under a restriction $\rho$ is built by
a fuel-limited recursion; write $G_k(f,\rho)$ for the decision tree it returns when run
with fuel budget $k \in \bbn$ starting from $\rho$, so that the canonical decision tree
of $f$ under $\rho$ is $G_{m(\rho)+1}(f,\rho)$, where $m(\rho)$ denotes the number of
free variables of $\rho$. Let $f$ be a DNF formula on $n$ variables, let $t$ be a term
with $t \in f$, and let $\rho_{\mathrm{orig}}$ and $\rho'$ be restrictions with
$m(\rho_{\mathrm{orig}}) \ge m(\rho') + 1$. Then the tree that returns the constant-true
leaf if $t$ is fixed by $\rho'$ and returns $G_{m(\rho_{\mathrm{orig}})}(f,\rho')$
otherwise is equal to the canonical decision tree of $f$ under $\rho'$.
\end{lemma}

\begin{lemma}[Unfolding the term sub-tree at a free head literal]
\lean{SwitchingLemma2.termSubTree_cons_free}
\label{SwitchingLemma2.termSubTree_cons_free}
\leanok
\uses{DecisionTree, Literal, SwitchingLemma2.Restriction, SwitchingLemma2.Restriction.freeVars, SwitchingLemma2.termSubTree}
\difficulty{1}
Let $l$ be a literal on $n$ Boolean variables, write $i$ for its underlying variable,
let $\mathit{rest}$ be a further list of literals, let $\rho$ be a restriction on these
$n$ variables, and let the continuation assign a decision tree to each restriction. If
$i$ is free under $\rho$, then the sub-tree built for the list $l :: \mathit{rest}$
under $\rho$ (with this continuation) is the branch that queries $i$, whose false-child
is the sub-tree built for $\mathit{rest}$ under the restriction obtained from $\rho$ by
fixing $i$ to false, and whose true-child is the sub-tree built for $\mathit{rest}$
under the restriction obtained from $\rho$ by fixing $i$ to true, both with the same
continuation.
\end{lemma}

\begin{lemma}[Skipping a fixed head literal in the term sub-tree]
\lean{SwitchingLemma2.termSubTree_cons_nonfree}
\label{SwitchingLemma2.termSubTree_cons_nonfree}
\leanok
\uses{DecisionTree, Literal, SwitchingLemma2.Restriction, SwitchingLemma2.Restriction.freeVars, SwitchingLemma2.termSubTree}
\difficulty{1}
Let $\rho$ be a restriction on $n$ variables, let $l$ be a literal whose underlying
variable $v$ is not free under $\rho$ (that is, $\rho$ already assigns $v$ a fixed
Boolean value), let $\mathit{rest}$ be a list of literals, and let $\mathrm{cont} :
\mathrm{Restriction}\,n \to \mathrm{DecisionTree}\,n$ be a continuation. Then the
sub-tree built for the term $l :: \mathit{rest}$ under $\rho$ with continuation
$\mathrm{cont}$ coincides with the one built for $\mathit{rest}$ under $\rho$ with the
same continuation; the fixed head literal $l$ contributes no branch.
\end{lemma}

\begin{lemma}[Head of the deepest path through a term sub-tree]
\lean{SwitchingLemma2.termSubTree_deepPath_head_free}
\label{SwitchingLemma2.termSubTree_deepPath_head_free}
\leanok
\uses{DecisionTree, DecisionTree.deepPath, Literal, SwitchingLemma2.Restriction, SwitchingLemma2.Restriction.freeVars, SwitchingLemma2.termSubTree, SwitchingLemma2.termSubTree_cons_free}
\difficulty{3}
Fix a restriction $\rho$ on $n$ variables and a continuation $c$ that assigns a decision
tree to each restriction. Let $l$ be a literal whose variable $v$ is free in $\rho$, and
let $\mathit{rest}$ be a further list of literals. Then there is a Boolean value $b$
such that the deepest root-to-leaf path of the term sub-tree built from the literal list
$l :: \mathit{rest}$, the restriction $\rho$, and the continuation $c$ begins with the
pair $(v, b)$ and continues with the deepest root-to-leaf path of the term sub-tree
built from $\mathit{rest}$, the restriction $\rho$ with $v$ now fixed to the value $b$,
and the same continuation $c$.
\end{lemma}

\begin{lemma}[Filtering literals by freeness is invariant under an unrelated update]
\lean{SwitchingLemma2.filter_free_update_eq}
\label{SwitchingLemma2.filter_free_update_eq}
\leanok
\uses{Literal, SwitchingLemma2.Restriction, SwitchingLemma2.Restriction.freeVars}
\difficulty{3}
Let $\rho : \mathrm{Fin}\,n \to \mathrm{Option}\,\mathrm{Bool}$ be a restriction, let
$R$ be a list of literals on the variables $\mathrm{Fin}\,n$, and let $v \in
\mathrm{Fin}\,n$ and $b \in \mathrm{Bool}$. Write $\rho[v \mapsto b]$ for the
restriction that agrees with $\rho$ everywhere except at $v$, where it takes the fixed
value $b$. If the variable of every literal in $R$ is distinct from $v$, then selecting
from $R$ those literals whose variable is free under $\rho[v \mapsto b]$ yields the same
list as selecting those whose variable is free under $\rho$.
\end{lemma}

\begin{lemma}[Deep-path variables track the free literals of a term]
\lean{SwitchingLemma2.termSubTree_deepPath_var_match}
\label{SwitchingLemma2.termSubTree_deepPath_var_match}
\leanok
\uses{DecisionTree, DecisionTree.deepPath, Literal, SwitchingLemma2.Restriction, SwitchingLemma2.Restriction.freeVars, SwitchingLemma2.filter_free_update_eq, SwitchingLemma2.termSubTree, SwitchingLemma2.termSubTree_cons_nonfree, SwitchingLemma2.termSubTree_deepPath_head_free}
\difficulty{5}
Let $\mathit{lits}$ be a list of literals on $n$ Boolean variables whose variables are
pairwise distinct, let $\rho$ be a restriction, and let $g$ be a map sending each
restriction to a decision tree. Let $T$ be the sub-tree built for the term
$\mathit{lits}$ under $\rho$ with continuation $g$: reading the literals of
$\mathit{lits}$ in order, it branches on each literal whose variable is free under
$\rho$, skips the others, and evaluates $g$ at each leaf. Write $\mathit{lits}_\rho$ for
the sublist of $\mathit{lits}$, in order, consisting of exactly those literals whose
variable is free under $\rho$. Then for every index $k$ smaller than both the length of
$\mathit{lits}_\rho$ and the length of the deepest root-to-leaf path of $T$, the
variable queried at position $k$ of that deepest path equals the variable of the $k$-th
literal of $\mathit{lits}_\rho$.
\end{lemma}

\begin{lemma}[Deepest-path length of a term sub-tree]
\lean{SwitchingLemma2.termSubTree_deepPath_append}
\label{SwitchingLemma2.termSubTree_deepPath_append}
\leanok
\uses{DecisionTree, DecisionTree.deepPath, Literal, SwitchingLemma2.Restriction, SwitchingLemma2.Restriction.freeVars, SwitchingLemma2.filter_free_update_eq, SwitchingLemma2.termSubTree, SwitchingLemma2.termSubTree_cons_nonfree, SwitchingLemma2.termSubTree_deepPath_head_free}
\difficulty{5}
Let $\rho$ be a restriction on $n$ variables, let $\mathit{lits}$ be a list of literals
whose underlying variables are pairwise distinct, and let $\mathrm{cont}$ be a
continuation assigning a decision tree to every restriction. Then there is a restriction
$\rho'$ that agrees with $\rho$ at every variable not occurring among the literals of
$\mathit{lits}$, and such that the deepest root-to-leaf path of the term sub-tree built
from $\mathit{lits}$, $\rho$, and $\mathrm{cont}$ has length equal to the number of
literals in $\mathit{lits}$ whose variable is free under $\rho$, plus the length of the
deepest root-to-leaf path of $\mathrm{cont}(\rho')$.
\end{lemma}

\begin{lemma}[Deepest-path split of a term's sub-tree]
\lean{SwitchingLemma2.termSubTree_deepPath_split}
\label{SwitchingLemma2.termSubTree_deepPath_split}
\leanok
\uses{DecisionTree, DecisionTree.deepPath, Literal, SwitchingLemma2.Restriction, SwitchingLemma2.Restriction.freeVars, SwitchingLemma2.filter_free_update_eq, SwitchingLemma2.termSubTree, SwitchingLemma2.termSubTree_cons_nonfree, SwitchingLemma2.termSubTree_deepPath_head_free}
\difficulty{5}
Let $\mathit{lits}$ be a list of literals on $n$ variables in which no two literals
share the same variable, let $\rho$ be a restriction on $n$ variables, and let $c$ be a
continuation, that is, a map sending each restriction to a decision tree. Then there
exist a list $P$ of (variable, branch-direction) pairs and a restriction $\rho'$ such
that the deepest path of the sub-tree for the term $\mathit{lits}$ built from $\rho$ and
$c$ equals $P$ followed by the deepest path of $c(\rho')$; the length of $P$ equals the
number of literals of $\mathit{lits}$ whose variable is free under $\rho$; and $\rho'$
agrees with $\rho$ on every variable that does not occur among the variables of
$\mathit{lits}$.
\end{lemma}

\begin{lemma}[Alive step of the canonical decision-tree builder]
\lean{SwitchingLemma2.canonicalDTree_alive_eq_termSubTree}
\label{SwitchingLemma2.canonicalDTree_alive_eq_termSubTree}
\leanok
\uses{DNF, SwitchingLemma2.Restriction, SwitchingLemma2.Term.fixedBy, SwitchingLemma2.Term.killedBy, SwitchingLemma2.termSubTree, Term}
\difficulty{2}
Fix $n$ Boolean variables, and let $f$ be a DNF formula (a list of terms) and $\rho$ a
restriction. Write $T(f, k, \rho)$ for the fuel-bounded canonical decision tree of $f$
under $\rho$ built with fuel budget $k$, and for a term $s$, a restriction $\sigma$, and
a continuation $c$ mapping restrictions to decision trees, write $S(s, \sigma, c)$ for
the term sub-tree of $s$ under $\sigma$ with continuation $c$. Suppose that, under
$\rho$, not every term of $f$ is killed and no term of $f$ is fixed, and let $t$ be the
first term of $f$, in list order, that is not killed by $\rho$. Then
\[
T(f,\ \mathit{fuel}+1,\ \rho) \;=\; S(t,\ \rho,\ c),
\]
where the continuation $c$ sends a restriction $\rho'$ to the leaf that outputs true
when $t$ is fixed by $\rho'$, and to $T(f,\ \mathit{fuel},\ \rho')$ otherwise.
\end{lemma}

\begin{lemma}[One-step unfolding of the canonical tree at a live restriction]
\lean{SwitchingLemma2.canonicalDTree_alive_eq_termSubTree'}
\label{SwitchingLemma2.canonicalDTree_alive_eq_termSubTree'}
\leanok
\uses{DNF, SwitchingLemma2.Restriction, SwitchingLemma2.Restriction.numFree, SwitchingLemma2.Term.fixedBy, SwitchingLemma2.Term.killedBy, SwitchingLemma2.canonicalDTree, SwitchingLemma2.canonicalDTree_alive_eq_termSubTree, SwitchingLemma2.termSubTree, Term}
\difficulty{1}
Let $f$ be a DNF formula and $\rho$ a restriction on $n$ variables such that $\rho$ does
not kill every term of $f$ and fixes no term of $f$, and let $t$ be the first term of
$f$, in list order, that $\rho$ does not kill. Then the canonical decision tree of $f$
under $\rho$ equals the term sub-tree built for $t$ over $\rho$ whose continuation,
evaluated at each leaf restriction $\rho'$ reached along a root-to-leaf path, returns
the constant-true leaf when $\rho'$ fixes $t$ and otherwise returns the fuel-bounded
canonical construction for $f$ under $\rho'$ run with fuel budget equal to the number of
free variables of $\rho$.
\end{lemma}

\begin{lemma}[Dropping a variable-fixed prefix from the term sub-tree]
\lean{SwitchingLemma2.termSubTree_skip_nonfree_prefix'}
\label{SwitchingLemma2.termSubTree_skip_nonfree_prefix'}
\leanok
\uses{DecisionTree, Literal, SwitchingLemma2.Restriction, SwitchingLemma2.Restriction.freeVars, SwitchingLemma2.termSubTree, SwitchingLemma2.termSubTree_cons_nonfree}
\difficulty{3}
Fix a restriction $\rho$ on $n$ variables and a continuation $c$ assigning a decision
tree to each restriction. For a list of literals $L$, write $T(L)$ for the term sub-tree
built from $L$ under $\rho$ with continuation $c$: it branches, in order, on each
literal of $L$ whose variable is free in $\rho$, skipping the rest, and at every leaf
invokes $c$ on the restriction updated along that path. Let $P$ and $R$ be lists of
literals, and suppose every literal in $P$ has its variable already fixed by $\rho$,
that is, not among the free variables of $\rho$. Then prepending $P$ to $R$ leaves the
resulting tree unchanged:
\[
T(P \frown R) = T(R),
\]
where $P \frown R$ denotes the concatenation of $P$ followed by $R$.
\end{lemma}

\begin{lemma}[Skipping a fixed head literal in the term sub-tree]
\lean{SwitchingLemma2.termSubTree_skip_updated_head}
\label{SwitchingLemma2.termSubTree_skip_updated_head}
\leanok
\uses{DecisionTree, Literal, SwitchingLemma2.Restriction, SwitchingLemma2.Restriction.freeVars, SwitchingLemma2.termSubTree, SwitchingLemma2.termSubTree_cons_nonfree}
\difficulty{2}
Let $\rho$ be a restriction on $n$ variables and $g$ a continuation assigning a decision
tree to each restriction on $n$ variables. Fix a variable $v$, a Boolean value $b$, and
a literal $\ell$ whose variable is $v$, and let $\mathit{rest}$ be a further list of
literals. Write $\rho[v \mapsto b]$ for the restriction that agrees with $\rho$
everywhere except at $v$, where it is fixed to the value $b$. Then the term sub-tree
built from the list obtained by prepending $\ell$ to $\mathit{rest}$, under $\rho[v
\mapsto b]$ with continuation $g$, coincides with the term sub-tree built from
$\mathit{rest}$ alone under the same restriction and the same continuation.
\end{lemma}

\begin{lemma}[Depth bounds decision-tree depth]
\lean{SwitchingLemma2.depth_ge_dtDepth}
\label{SwitchingLemma2.depth_ge_dtDepth}
\leanok
\proofsource{switching-lemma}{pdf-b5e074215b9e-p001-b005}
\uses{DecisionTree, DecisionTree.depth, DecisionTree.eval, dtDepth}
\difficulty{1}
Let $f : (\mathrm{Fin}\,n \to \mathrm{Bool}) \to \mathrm{Bool}$ be a Boolean function,
and let $T$ be a decision tree on $n$ Boolean variables that computes $f$, in the sense
that $T$ evaluates to $f(x)$ on every input $x$. Then the depth of $T$ is at least the
decision-tree depth of $f$; that is, $\mathrm{depth}(T)$ is no smaller than the least
depth of any decision tree computing $f$.
\end{lemma}

\begin{lemma}[Canonical tree depth lower-bounds decision-tree depth]
\lean{SwitchingLemma2.canonicalDTree_depth_ge}
\label{SwitchingLemma2.canonicalDTree_depth_ge}
\leanok
\uses{DNF, DNF.eval, DecisionTree.depth, SwitchingLemma2.Restriction, SwitchingLemma2.canonicalDTree, SwitchingLemma2.canonicalDTree_correct, SwitchingLemma2.depth_ge_dtDepth, SwitchingLemma2.restrictFn, dtDepth}
\difficulty{1}
Let $f$ be a DNF formula on $n$ Boolean variables and let $\rho$ be a restriction on
those variables, and write $f|_\rho$ for the restricted Boolean function $x \mapsto
f\big(\rho.\mathrm{extend}\,x\big)$, which evaluates $f$ on the assignment that takes
the fixed value of $\rho$ where one is set and reads the remaining coordinates from $x$.
Then the depth of the canonical decision tree for $f$ under $\rho$ is at least the
decision-tree depth of $f|_\rho$; that is, the least depth of any decision tree
computing $f|_\rho$ does not exceed the depth of the canonical tree.
\end{lemma}

\begin{lemma}[Decision-tree depth bounded by the number of free variables]
\lean{SwitchingLemma2.dtDepth_restrictFn_le_numFree}
\label{SwitchingLemma2.dtDepth_restrictFn_le_numFree}
\leanok
\uses{DecisionTree, DecisionTree.depth, DecisionTree.eval, SwitchingLemma2.Restriction, SwitchingLemma2.Restriction.extend, SwitchingLemma2.Restriction.freeVars, SwitchingLemma2.Restriction.numFree, SwitchingLemma2.depth_ge_dtDepth, SwitchingLemma2.extend_update_self, SwitchingLemma2.numFree_update_lt, SwitchingLemma2.restrictFn, dtDepth}
\difficulty{5}
Let $f : (\mathrm{Fin}\,n \to \mathrm{Bool}) \to \mathrm{Bool}$ be a Boolean function
and let $\rho$ be a restriction on $n$ variables, fixing some coordinates to Boolean
values and leaving the rest free. Then the decision-tree depth of the restricted
function $f\!\restriction_\rho$, which sends $x$ to $f$ evaluated with the free
coordinates taken from $x$ and the remaining coordinates set by $\rho$, is at most the
number of free variables of $\rho$.
\end{lemma}
